\section{Conclusions}
\label{sec:conclusions}

\begin{enumerate}[leftmargin=1.6em,itemsep=4pt,topsep=4pt,parsep=0pt]

\item No published radio technosignature search above about 30\,GHz has
reached a star within 40\,pc (\S\ref{sec:background}), and an archive
assembled for other science is at present the only way to reach one.

\item The primary experiment is \NWinA{} fine-channel windows toward
\NSysClassA{} stellar systems, covering \CvUnionA\,GHz of union bandwidth in
disconnected intervals, searched for an unresolved carrier whose frequency
drifts. A supplementary coarse-channel search for an unresolved spectral
excess adds \NWinB{} windows toward \PrSysClassB{} systems and
\CvUnionB\,GHz of union bandwidth, with no drift discrimination. Both come
from one uniform reduction of \NEB{} public execution blocks,
\TotalDataTB\,TB of raw data, and between them they span
\SurvFreqLo--\SurvFreqHi\,GHz.

\item Sensitivity is a measured completeness, not a threshold:
$\mathrm{EIRP}_{90}$ is the power at which nine of ten carriers injected into
the real calibrated visibilities come back through the unchanged selection,
the trigger and localisation at the star. The median
system's best window reaches $\EirpNinetySysMedA$\,W, with a spread across
systems of $\EirpNinetySysLoA$--$\EirpNinetySysHiA$\,W. The rank and the
visibility fit decide nothing, so neither is charged to that limit. One bias
runs one way only and is not inside the quoted calibration: an injected
carrier suffers no atmospheric decorrelation where a real one would, and the
coherence each block actually lost, measured on its own phase calibrator,
makes the median limit optimistic by \SensDecorPct{} per cent and those above
\SensDecorHiGHz{}\,GHz by \SensDecorHiPct{} per cent.

\item All \EvNCrossRaw{} threshold crossings are reported with their
dispositions, \PrFlagXrossWord{} of them in the one execution block
that fails the data-quality criterion of
Appendix~\ref{sec:blockquality} and reported apart from the
\EvNCrossQp{} in quality-passing data; of those, \EvNAttrQp{} coincide with
a catalogued molecular transition, almost all of them carbon monoxide ---
including the $\beta$~Pictoris disc, which the search recovered without
being told where to look --- and \EvNUnattrQp{} do not. Of the unattributed,
\MkNSODecideWord{} coincide with sulphur monoxide, which no debris disc is known to
carry and which here would have to outshine an undetected carbon monoxide
line in the same window; the coincidence is reported and not acted on, and
disposing of them instead would make the split over all \EvNCrossRaw{}
crossings \MkAttrAlt{} to \MkUnattrAlt{} rather than \MkAttr{} to
\MkUnattr. Chance, measured from each window's own control
positions and channel grid, predicts \CeExp{} unattributed Class~A
crossings with \CeLo--\CeHi{} allowed by resampling the execution blocks,
so the count is consistent with chance within a range that would not
register a small added population either way.

\item Every unattributed crossing was tested for recurrence in the star's
own frame, in each further execution block covering the event frequency and
over the whole of each such window at every drift searched, because an
accelerating carrier need not return to one channel. None was found
again: the largest statistic at any predicted cell is \RcTMax{}
against $T_{\rm trig}=\StageNullTrig$, and the exceedances found anywhere in
those windows are no more than chance predicts for trials of that size. Nor
did \PrRankClause{} return. Of the \RcNExcl{} crossings that retain a
discovery amplitude and so carry an exclusion, the repeat coverage rules out
a carrier at a median \RcFracMed{} of that amplitude, with the individual
limits running from \RcFracMin{} to \RcFracMax: the median is not a bound on
the set. Where the repeat is shallow the criterion has little power and its
silence is not an exclusion --- a carrier radiating at the discovery
amplitude through both observations would have met it with probability above
\RcPowerLevel{} per cent for \RcNPowerHi{} of those crossings but only
\RcPowerMin--\RcPowerLoHi{} for the remaining \RcNPowerLo, which are
untested rather than excluded. What these tests exclude is a carrier that
holds its frequency inside the window searched and radiates through both
observations, and not a persistent transmitter of any kind.

\item What this survey could search and what it could confirm are different
extents: of the \RcSearchSys{} systems the carrier search covered, only
\RcConfSysDay{} were observed again at the same frequency more than a day
apart and \RcConfSysYr{} beyond a year. Across both classes,
\RcNoRepSysAll{} of the \RcSearchSysAll{} systems have no second block
covering a frequency the first one covered, and toward those the null
reported here says nothing at any strength.

\item Combining each surveyed star's epochs in its own rest frame deepens
the flux limit by a median factor \SyStkGainMed. For Proxima~Centauri,
where $\sqrt{N_{\rm eff}}=\SxProxSqrtNeff$ is what combining the epochs
buys, the stacked limit measured by injection into the stack reaches
$\SzProxEirpCensus$\,W over the primary census alone and
$\StkProxEirpCal$\,W once the rest of the archive is added, in
\ChanBModeMHz\,MHz channels --- a Class~B limit on an unresolved excess,
carrying no drift discrimination --- for an emitter steady in the stellar
frame to within one channel, \StkProxTolKms\,km\,s$^{-1}$ peak to peak,
across \StkProxSpanYr\,yr.

\item These results constrain unresolved carriers only within the domain
actually searched and only toward this archive-selected sample
(\S\ref{sec:excludes}); what transfers beyond it is the obstacle, that a
channelisation chosen by other observers for other purposes, and not
collecting area, decides which kind of search is possible at all.

\end{enumerate}
